\documentclass[UTF8]{article}
\usepackage[UTF8]{ctex}
\usepackage{amsmath, amssymb}
\usepackage{physics}
\usepackage{graphicx}
\usepackage{subcaption}
\usepackage{caption}
\usepackage[font=small]{caption}
\usepackage{booktabs}
\usepackage{multirow}
\usepackage{geometry}
\geometry{a4paper, left=2.5cm, right=2.5cm, top=2.5cm, bottom=2.5cm}
\usepackage{enumitem}
\usepackage{float}

% 页眉页脚
\usepackage{fancyhdr}
\pagestyle{fancy}
\lhead{色散与偏振实验}
\chead{}
\rhead{物理实验B}
\lfoot{}
\cfoot{\thepage}
\rfoot{}
\renewcommand{\headrulewidth}{0.4pt}
\setlength{\headheight}{14pt}

\title{色散与偏振实验}
\author{} % 填写你的姓名和学号
\date{} % 填写实验日期

\begin{document}

\maketitle

% ============================================================
\section{实验目的}
% ============================================================

\begin{enumerate}
    \item 了解光的色散现象，理解折射率随波长变化的规律。
    \item 学习搭建分离式平台分光光谱仪，掌握\textbf{最小偏向角法}测定三棱镜折射率的原理与方法。
    \item 学习利用\textbf{布儒斯特角法}测定介质折射率，理解偏振光的产生机制。
    \item 观察玻璃球和亚克力棒中的\textbf{虹}与\textbf{霓}现象，验证其偏振特性。
    \item 学习使用\textbf{光栅光谱仪}测量和分析不同光源的发射光谱。
    \item 观察\textbf{特斯拉线圈}的放电现象，理解无线发光和电弧放电的原理。
\end{enumerate}

% ============================================================
\section{实验原理}
% ============================================================

\subsection{光的色散}

当白光通过三棱镜时，不同波长的光因折射率不同而发生不同程度的偏折，形成光谱，这一现象称为\textbf{色散}。对于绝大多数透明介质，折射率$n$随波长$\lambda$的减小而增大，即
\[
\frac{dn}{d\lambda} < 0
\]
这称为\textbf{正常色散}。柯西（Cauchy）色散公式可以近似描述这种关系：
\[
n(\lambda) = A + \frac{B}{\lambda^2} + \frac{C}{\lambda^4} + \cdots
\]
其中$A$、$B$、$C$为与材料相关的常数。

介质的色散能力通常用\textbf{阿贝数}（Abbe number）来衡量：
\[
\nu = \frac{n_D - 1}{n_F - n_C}
\]
其中$n_D$、$n_F$、$n_C$分别为钠黄光（589.3\,nm）、氢蓝光（486.1\,nm）和氢红光（656.3\,nm）对应的折射率。阿贝数越小，色散越强。

\subsection{最小偏向角法测折射率}

\subsubsection{偏向角}

光线通过三棱镜时，入射光和出射光之间的夹角称为\textbf{偏向角}~$\delta$。对于顶角为$A$的三棱镜，偏向角与入射角$\theta_1$、出射角$\theta_2$的关系为：
\[
\delta = \theta_1 + \theta_2 - A
\]

\subsubsection{最小偏向角条件}

当入射角等于出射角（$\theta_1 = \theta_2$）时，光线在棱镜内对称传播，此时偏向角取得最小值$\delta_{\min}$。这意味着棱镜内部的光线平行于底面。

在最小偏向角条件下，有：
\[
\theta_1 = \frac{A + \delta_{\min}}{2}, \qquad r = \frac{A}{2}
\]
其中$r$为棱镜内的折射角。由折射定律$n = \sin\theta_1 / \sin r$可得：
\begin{equation}
\boxed{n = \frac{\sin\dfrac{A + \delta_{\min}}{2}}{\sin\dfrac{A}{2}}}
\label{eq:n-min-dev}
\end{equation}

对于等边三棱镜，$A = 60^\circ$，上式简化为：
\[
n = \frac{\sin\dfrac{60^\circ + \delta_{\min}}{2}}{\sin 30^\circ} = 2\sin\dfrac{60^\circ + \delta_{\min}}{2}
\]

\subsubsection{实验中寻找最小偏向角}

实验中通过缓慢转动载有三棱镜的转台，观察望远镜中谱线的运动方向。当谱线即将反向运动的临界位置，即为该谱线对应的\textbf{最小偏向角位置}。此时锁定转台，精确读取出射方位角$\varphi$；取下三棱镜后，将望远镜对准直射狭缝光，读取入射方位角$\varphi_0$。偏向角为：
\[
\delta_{\min} = \varphi - \varphi_0
\]

\subsection{布儒斯特角与偏振}

\subsubsection{布儒斯特定律}

当自然光以特定角度$\theta_B$（布儒斯特角）入射到两种介质的界面时，反射光成为完全的\textbf{线偏振光}（仅含S分量），反射光与折射光互相垂直。此时折射率$n$与布儒斯特角$\theta_B$满足：
\begin{equation}
\boxed{n = \tan\theta_B}
\label{eq:brewster}
\end{equation}

\subsubsection{物理解释}

当入射角为布儒斯特角时，折射光的传播方向与反射光的传播方向恰好垂直（$\theta_B + \theta_r = 90^\circ$）。由于光是横波，折射光方向上振动的P分量恰好沿反射方向，无法激发反射波中的P分量，因此反射光中P分量为零，仅剩S分量。

\subsubsection{实验测定方法}

用半导体激光器照射三棱镜的光学平面，旋转偏振片使入射光成为P偏振光（振动方向在入射面内），然后旋转转台调节入射角，寻找反射光\textbf{最暗}的位置。该位置对应的入射角即为布儒斯特角。

\subsection{虹与霓}

\subsubsection{虹（主虹）}

白光射入球形水滴（或实验中的玻璃球）后，经过\textbf{一次全内反射}从水滴折射出来。不同颜色的光因折射率不同而以不同角度出射，形成虹。虹的出射角约为$42^\circ$（对水），且\textbf{外侧为红色、内侧为紫色}。

\subsubsection{霓（副虹）}

光线在水滴内经过\textbf{两次全内反射}后出射形成霓。霓位于虹的外侧，颜色排列与虹\textbf{相反}（外紫内红），且亮度较虹更弱。虹和霓之间的暗带称为\textbf{亚历山大暗带}。

\subsubsection{偏振特性}

虹和霓的光具有偏振特性，其中\textbf{S偏振光（振动方向沿彩虹圆环切线）}占主导地位。可用偏振片旋转验证。

\subsection{光栅光谱仪}

光栅光谱仪利用\textbf{光栅衍射}原理分析光源的光谱成分。光栅的衍射方程为：
\[
d \sin\theta = m\lambda \qquad (m = 0, \pm 1, \pm 2, \ldots)
\]
其中$d$为光栅常数，$\theta$为衍射角，$m$为衍射级次。

光谱仪的\textbf{分辨本领}$R$定义为：
\[
R = \frac{\lambda}{\Delta\lambda} = mN
\]
其中$N$为光栅总缝数。分辨本领越高，能分辨的波长差越小。

\subsection{特斯拉线圈}

特斯拉线圈（Tesla Coil）是一种谐振变压器，基本原理是利用LC振荡回路产生高频高压电场。主要由初级线圈、次级线圈和放电端三部分组成。初级线圈与电容构成LC谐振回路，通过电磁感应在次级线圈产生极高的电压（可达数十万伏），从而电离空气产生电弧放电。

特斯拉线圈可以无线点亮灯泡和灯管，其机理分别为：LED灯泡通过电磁感应产生感应电流；气体放电管（节能灯管）因高频电场直接电离管内气体而发光；不同气体放电管因气体种类不同而呈现不同颜色。

% ============================================================
\section{实验仪器与设备}
% ============================================================

\subsection{色散实验}
\begin{itemize}
    \item 氦放电管（高压5$\sim$6\,kV，\textbf{严禁触碰电极}）
    \item 等边三棱镜（顶角$A = 60^\circ$，重火石玻璃 ZF）
    \item 可调狭缝、凸透镜（焦距$f = 20$\,cm）
    \item 望远镜（带分划板十字丝）
    \item 分离式转台（带角度刻度盘，估读精度$0.1^\circ$）
\end{itemize}

\subsection{偏振实验}
\begin{itemize}
    \item 半导体绿光激光器
    \item 偏振片（$\times$2）
    \item 玻璃球（大、小各一）
    \item 亚克力棒（多种直径）
    \item 白色大屏
\end{itemize}

\subsection{光谱仪实验}
\begin{itemize}
    \item 光栅光谱仪（USB接口）
    \item 电脑及光谱分析软件（窗台小圃光谱仪软件）
    \item 氦放电管、氢放电管、荧光灯等多种光源
\end{itemize}

\subsection{特斯拉线圈}
\begin{itemize}
    \item 特斯拉线圈装置
    \item LED灯泡、节能灯管、8种气体放电管
\end{itemize}

% ============================================================
\section{实验内容与步骤}
% ============================================================

\subsection{实验一：最小偏向角法测三棱镜折射率}

\begin{enumerate}
    \item \textbf{等高共轴粗调}：按照"氦灯$\to$狭缝$\to$凸透镜$\to$转台$\to$望远镜"顺序摆放各器件，目测调至等高。
    \item \textbf{调焦}：氦灯靠近狭缝，调节凸透镜距狭缝约20\,cm，在望远镜中观察到清晰无视差的狭缝像。调整狭缝宽度适中，十字丝对准清晰。
    \item \textbf{寻找光谱}：将三棱镜置于转台上（折射面朝向入射光），先用肉眼在出射侧搜索彩色光谱，再将望远镜对准该方向。
    \item \textbf{测量最小偏向角}：缓慢转动转台使谱线偏向角减小，至谱线即将反向运动的临界位置即为最小偏向角。锁死转台，用望远镜分划板对准目标谱线，记录出射方位角$\varphi$。
    \item \textbf{测入射方位角}：取下三棱镜，望远镜对准狭缝像，记录入射方位角$\varphi_0$。
    \item 对紫色（447.1\,nm）和黄色（587.6\,nm）两条谱线分别测量。
\end{enumerate}

\subsection{实验二：布儒斯特角法测折射率}

\begin{enumerate}
    \item 三棱镜置于转台上，用激光器照射其光学平面，调节使反射光与入射光重合（确认法线方向），记录法线方位角。
    \item 在激光器与转台之间放置偏振片，旋转偏振片使入射光为P偏振光。
    \item 旋转转台改变入射角（预设$\theta_i \approx 57^\circ$），同时观察白屏上反射光斑亮度。
    \item 反复微调入射角和偏振片方位，使反射光斑\textbf{最暗}，读取此时方位角。
    \item 布儒斯特角$\theta_B$等于当前方位角与法线方位角之差。
\end{enumerate}

\subsection{实验三：虹与霓的观察}

\begin{enumerate}
    \item \textbf{白光观察}：将大/小玻璃球依次放在转台上，用手电筒照射，观察一次反射形成的虹和二次反射形成的霓。
    \item \textbf{偏振验证}：用布儒斯特角实验中已校正的偏振片旋转观察，验证虹和霓中S光占主导。
    \item \textbf{激光观察}：用激光器照射亚克力棒，观察幻日、虹和霓现象，测定出射角度。
\end{enumerate}

\subsection{实验四：光谱仪测量}

\begin{enumerate}
    \item 光谱仪通过USB连接电脑，启动光谱分析软件。
    \item 将氦放电管对准光谱仪通光孔，调节高度和角度使信号最强。
    \item 调节软件中的选区框高度和位置，手动设置曝光强度使主峰尖锐。
    \item \textbf{波长校准}：使用氦光谱的447.1\,nm和706.6\,nm两个已知峰进行手动校准。
    \item 记录其他各峰的波长值，与标准氦光谱对比。
    \item 更换光源（氢管、荧光灯等），重复测量并记录光谱图。
\end{enumerate}

\subsection{实验五：特斯拉线圈观察}

\begin{enumerate}
    \item 了解特斯拉线圈的基本构造（LC振荡回路、升压线圈、放电端）。
    \item 通电观察电弧放电现象。
    \item 分别用特斯拉线圈无线点亮LED灯泡、节能灯管和8种不同气体的放电管，观察并比较发光颜色和亮度。
    \item 利用特斯拉线圈播放音乐，分析电弧放电产生声音的物理机制。
\end{enumerate}

% ============================================================
\section{数据处理与结果}
% ============================================================

\subsection{实验一：最小偏向角法}

\subsubsection{实验数据}

等边三棱镜顶角 $A = 60^\circ$。

\begin{table}[H]
\centering
\caption{最小偏向角法——原始数据}
\begin{tabular}{cccc}
\toprule
谱线 & 波长 (nm) & 出射角 $\varphi$ & 入射角 $\varphi_0$ \\
\midrule
紫色 & 447.1 & $295.4^\circ$ & $242.0^\circ$ \\
黄色 & 587.6 & $258.8^\circ$ & $205.5^\circ$ \\
\bottomrule
\end{tabular}
\end{table}

\subsubsection{最小偏向角计算}

\[
\delta_{\min} = \varphi - \varphi_0
\]

紫色：$\delta_{\text{紫}} = 295.4^\circ - 242.0^\circ = 53.4^\circ$

黄色：$\delta_{\text{黄}} = 258.8^\circ - 205.5^\circ = 53.3^\circ$

\subsubsection{折射率计算}

由式~(\ref{eq:n-min-dev})，代入$A = 60^\circ$：

\paragraph{紫色谱线（$\delta = 53.4^\circ$）}

\[
n_{\text{紫}} = \frac{\sin\dfrac{60^\circ + 53.4^\circ}{2}}{\sin 30^\circ}
= \frac{\sin 56.70^\circ}{0.5}
= \frac{0.8358}{0.5}
= \boxed{1.672}
\]

\paragraph{黄色谱线（$\delta = 53.3^\circ$）}

\[
n_{\text{黄}} = \frac{\sin\dfrac{60^\circ + 53.3^\circ}{2}}{\sin 30^\circ}
= \frac{\sin 56.65^\circ}{0.5}
= \frac{0.8353}{0.5}
= \boxed{1.671}
\]

\subsubsection{不确定度计算（以黄色谱线为例）}

\paragraph{偏向角的不确定度}

角度估读精度 $\Delta\varphi = 0.1^\circ$，则：
\[
\Delta\delta = \sqrt{(\Delta\varphi)^2 + (\Delta\varphi_0)^2}
= \sqrt{0.1^2 + 0.1^2}
= 0.1\sqrt{2}
\approx 0.141^\circ
\]

转为弧度：
\[
\Delta\delta = 0.141^\circ \times \frac{\pi}{180} = 0.00246 \text{ rad}
\]

\paragraph{折射率的不确定度}

\[
\frac{\partial n}{\partial \delta}
= \frac{\cos\dfrac{A+\delta}{2}}{2\sin\dfrac{A}{2}}
= \frac{\cos 56.65^\circ}{2 \times 0.5}
= \frac{0.5497}{1}
= 0.5497
\]

\[
\Delta n = \left|\frac{\partial n}{\partial \delta}\right| \times \Delta\delta
= 0.5497 \times 0.00246
\approx \boxed{0.001}
\]

\paragraph{最终结果}

\[
n_{\text{黄}} = 1.671 \pm 0.001
\]

不确定度 $\Delta n = 0.001$ 影响到小数点后第三位，因此 $n$ 保留四位有效数字。

\subsection{实验二：布儒斯特角法}

\subsubsection{布儒斯特角测定}

采用绿光半导体激光器照射三棱镜光学面。旋转转台和偏振片，找到反射光斑最暗的位置。实验数据如下：

\begin{itemize}
    \item 入射面法向所在位置标度：$58.6^\circ$
    \item 布儒斯特角位置标度：$116.0^\circ$
\end{itemize}

布儒斯特角：
\[
\theta_B = 116.0^\circ - 58.6^\circ = 57.4^\circ
\]

\subsubsection{折射率计算}

由式~(\ref{eq:brewster})：
\[
n = \tan \theta_B = \tan 57.4^\circ \approx \boxed{1.564}
\]

\subsection{结果汇总}

\begin{table}[H]
\centering
\caption{最小偏向角法——计算结果汇总}
\begin{tabular}{ccccccc}
\toprule
谱线 & 波长 & 出射角 $\varphi$ & 入射角 $\varphi_0$ & 最小偏向角 $\delta$ & 折射率 $n$ & 不确定度 $\Delta n$ \\
     & (nm)  &                  &                    &                     &            &                      \\
\midrule
紫色 & 447.1 & $295.4^\circ$ & $242.0^\circ$ & $53.4^\circ$ & 1.672 & — \\
黄色 & 587.6 & $258.8^\circ$ & $205.5^\circ$ & $53.3^\circ$ & 1.671 & $\pm 0.001$ \\
\bottomrule
\end{tabular}
\end{table}

\begin{table}[H]
\centering
\caption{布儒斯特角法——计算结果}
\begin{tabular}{ccccc}
\toprule
光源 & 法线方位 & 布儒斯特角方位 & 布儒斯特角 $\theta_B$ & 折射率 $n = \tan\theta_B$ \\
\midrule
绿光激光 & $58.6^\circ$ & $116.0^\circ$ & $57.4^\circ$ & 1.564 \\
\bottomrule
\end{tabular}
\end{table}

\begin{table}[H]
\centering
\caption{不同方法测定折射率汇总}
\begin{tabular}{cccc}
\toprule
测量方法 & 测量对象 & 折射率 $n$ & 备注 \\
\midrule
最小偏向角法 & 紫色谱线 (447.1\,nm) & 1.672 & 重火石玻璃棱镜 \\
最小偏向角法 & 黄色谱线 (587.6\,nm) & $1.671 \pm 0.001$ & 重火石玻璃棱镜 \\
布儒斯特角法 & 绿光激光            & 1.564 & 三棱镜光学面 \\
\bottomrule
\end{tabular}
\end{table}

% ============================================================
\section{分析与讨论}
% ============================================================

\subsection{两种方法测定折射率的比较}

最小偏向角法测得的三棱镜折射率约为1.67，而布儒斯特角法测得的折射率为1.564，两者存在明显差异。这一差异可从以下几个方面解释：

\begin{enumerate}
    \item \textbf{测量对象不同}：最小偏向角法测定的是三棱镜棱镜体材料（重火石玻璃 ZF）的折射率，而布儒斯特角法测定的是三棱镜光学面的折射率。若光学面上有镀膜或表面处理层，则测定结果可能反映的是表面层的折射率，而非体材料的折射率。
    
    \item \textbf{色散效应}：最小偏向角法使用了氦灯的特征谱线（447.1\,nm和587.6\,nm），而布儒斯特角法使用绿光激光（约532\,nm）。由色散关系$n(\lambda)$，不同波长的折射率本身就不同。但即使考虑色散效应，紫色和黄色谱线的折射率差异仅为0.001，远不足以解释$1.67$与$1.56$之间约0.1的差距。
    
    \item \textbf{测量精度差异}：最小偏向角法利用了对称光路条件，系统误差较小，测量精度较高（不确定度$\Delta n \approx 0.001$）。布儒斯特角法依赖于目视判断反射光最暗位置，主观性较强，且布儒斯特角附近P分量反射率变化较为平缓，精确定位困难。
    
    \item \textbf{最可能的原因}：重火石玻璃（ZF系列）的折射率范围通常在1.6$\sim$1.8之间，最小偏向角法测得的$n \approx 1.67$与标准值吻合良好。布儒斯特角法测得的$n \approx 1.56$更接近冕牌玻璃（K系列）的折射率范围，暗示所测光学面可能并非棱镜体材料本身，或存在系统误差。
\end{enumerate}

\subsection{色散现象的验证}

由实验数据可知，紫色谱线（447.1\,nm）对应的折射率$n_{\text{紫}} = 1.672$略大于黄色谱线（587.6\,nm）的$n_{\text{黄}} = 1.671$，验证了\textbf{正常色散}规律：短波长（高频率）的光在介质中折射率更大。

然而，两条谱线的折射率差异$\Delta n = 0.001$恰好与不确定度的量级相当。若需更精确地研究色散关系，应选取更多条谱线（氦灯具有丰富的可见光谱线：紫色447.1\,nm、蓝绿色501.6\,nm、黄色587.6\,nm、橙红色667.8\,nm、红色706.5\,nm等），绘制$n$-$\lambda$曲线，拟合柯西色散方程。

\subsection{虹和霓的偏振特性}

实验中通过旋转偏振片观察虹和霓，验证了其光的偏振特性。虹和霓均以S偏振光为主导，偏振方向沿彩虹圆环的切线方向。这与菲涅尔反射公式的理论预测一致：在接近布儒斯特角的入射条件下，P分量的反射系数远小于S分量。

\subsection{误差来源分析}

\begin{enumerate}
    \item \textbf{角度读数误差}：分度盘的最小刻度为$1^\circ$，估读至$0.1^\circ$，是折射率不确定度的主要来源。
    \item \textbf{最小偏向角的确定误差}：谱线即将反向运动的临界位置不易精确判断，存在一定的主观误差。
    \item \textbf{等高共轴调节}：若光路未严格等高共轴，测量的偏向角将偏离理想值。
    \item \textbf{狭缝宽度}：狭缝过宽会导致谱线展宽，难以精确对准十字丝。
    \item \textbf{三棱镜顶角}：假设$A = 60^\circ$精确成立，实际三棱镜加工可能存在微小偏差。
\end{enumerate}

\subsection{特斯拉线圈发声原理}

特斯拉线圈产生的电弧放电会加热电弧通道中的空气，使其剧烈膨胀形成声波。通过调制驱动信号的频率，使电弧以音频频率脉动放电，便可产生对应频率的声波，实现"播放音乐"的效果。这本质上是将电能直接转化为声能的等离子体扬声器原理。

% ============================================================
\section{思考与总结}
% ============================================================

本实验通过多种方法研究了光与物质相互作用中的色散和偏振现象。最小偏向角法提供了一种高精度测量折射率的手段，其不确定度可达$\pm 0.001$量级；布儒斯特角法虽精度较低，但直观地展示了偏振光的产生机制。虹和霓的观察将几何光学和波动光学的概念有机结合，加深了对全内反射和色散的理解。光谱仪实验则将传统分光方法与现代光电检测技术进行了对比，展示了技术进步对实验物理的推动。

% ============================================================
% \section{参考文献}
% ============================================================
% 如需添加参考文献，可在此处补充

\end{document}
